\subsubsection{মডেল, সহায়ক সফটওয়্যার ও ব্যবহার}
মডেল ডেটা থেকে শেখা গাণিতিক সম্পর্ক ও প্যারামিটার ব্যবহার করে ইনপুটের ফল নির্ণয় করে। একই ধারার একাধিক মডেলকে \term{Model Family} বলে; তাদের আকার, গতি ও ক্ষমতা ভিন্ন হতে পারে। \term{App বা Assistant} হলো মডেল ব্যবহার করার সফটওয়্যার। তাতে কথোপকথন, ফাইল, অনুসন্ধান ও অন্য টুল যুক্ত থাকতে পারে। অ্যাপের সুবিধা প্রতিটি মডেলের নিজস্ব ক্ষমতা নয়।
ChatGPT হলো OpenAI-এর কথোপকথনভিত্তিক সহায়ক সফটওয়্যার; GPT পরিচিত মডেল পরিবার। Claude, Grok ও Gemini নাম মডেল পরিবার এবং সহায়ক সফটওয়্যার---দুই প্রসঙ্গেই ব্যবহৃত হয়। তাই কোন মডেল ও কোন সফটওয়্যারের সুবিধা ব্যবহার হচ্ছে তা আলাদা করে বোঝা দরকার।
\begin{compacttable}
\begin{longtable}{|p{\dimexpr.19\linewidth-7pt\relax}|p{\dimexpr.24\linewidth-7pt\relax}|p{\dimexpr.57\linewidth-7pt\relax}|}
\hline\textbf{পরিবার} & \textbf{নির্মাতা} & \textbf{কাজের পরিচিত ধরন}\\\hline\endfirsthead
\hline\textbf{পরিবার} & \textbf{নির্মাতা} & \textbf{কাজের পরিচিত ধরন}\\\hline\endhead
GPT & OpenAI & ভাষা, বিশ্লেষণ ও কোড; নির্বাচিত মডেলে ছবি বোঝা ও টুল ব্যবহারের সমর্থন। HTML-এর ভুল ব্যাখ্যা বা তথ্য থেকে প্রতিবেদন তৈরি করা যায়।\\\hline
Claude & Anthropic & লেখা, নথি বিশ্লেষণ ও কোড; নির্বাচিত মডেলে ছবি ও টুল। দীর্ঘ নোট তুলনা ও কোড পরীক্ষা করা যায়।\\\hline
Grok & xAI & ভাষা, যুক্তিভিত্তিক কাজ ও কোড; নির্বাচিত মডেলে দৃষ্টি ও টুল। অনুসন্ধান যুক্ত থাকলে নতুন তথ্যের উৎস খোঁজা যায়।\\\hline
Gemini & Google & নির্বাচিত মডেলে লেখা, ছবি, শব্দ ও ভিডিও বিশ্লেষণ। দৃশ্য ও কথার সম্পর্ক থেকে সারসংক্ষেপ তৈরি করা যায়।\\\hline
DeepSeek & DeepSeek & ভাষা, যুক্তি ও কোড। প্রকাশিত মডেল এবং সেবার সুবিধা আলাদা হতে পারে।\\\hline
Llama & Meta & প্রকাশিত ওজনসহ মডেল পরিবার। নির্দিষ্ট লাইসেন্স ও উপযুক্ত যন্ত্রে নিজস্ব সহায়ক ব্যবস্থা তৈরি করা যায়।\\\hline
Qwen & Alibaba-এর Qwen Team & ভাষা, কোড ও নির্বাচিত বহুমাধ্যম মডেল। প্রকাশিত ওজন ও সেবাভিত্তিক সংস্করণ আছে।\\\hline
Mistral & Mistral AI & সাধারণ ও বিশেষায়িত মডেল; কোড, নথি/OCR ও অডিওর পৃথক মডেলও আছে।\\\hline
\end{longtable}
\end{compacttable}
Google-এর কথোপকথনভিত্তিক সহায়ক Bard-এর নাম 2024 সালে Gemini রাখা হয়। মডেল পরিবার ও সেই মডেল ব্যবহারকারী সফটওয়্যারের পরিচয় এখানে পৃথক।
Microsoft Copilot-ও AI সহায়কের উদাহরণ। লেখা, সারসংক্ষেপ ও তথ্য খোঁজায় শেখানো মডেল ও ভাষা প্রক্রিয়াকরণ ব্যবহার হয়। Azure হলো Microsoft-এর ক্লাউড প্ল্যাটফর্ম; সহায়ক সফটওয়্যার, ক্লাউড সেবা ও AI-এর শাখা পৃথক পরিচয়।
একই পরিবারের সব মডেল একই ইনপুট বা আউটপুট সমর্থন করে না। ছবি বোঝা ও তৈরি, শব্দ বোঝা ও তৈরি পৃথক ক্ষমতা। অনুসন্ধান টুল সত্যিই ব্যবহৃত হলে নতুন তথ্য পাওয়া যেতে পারে; শুধু নাম থেকে সরাসরি অনুসন্ধান হয়েছে বলা যায় না। Web বা X থেকে উদ্ধার করা তথ্যেও তারিখ, মূল উৎস ও অসম্পূর্ণতা পরীক্ষা দরকার।
\subtopichead{কাজ অনুযায়ী মডেল}
\begin{compactpoints}
\item ভাষা মডেল লেখা ও প্রসঙ্গ থেকে ব্যাখ্যা, অনুবাদ, কোড বা নির্দিষ্ট গঠনের উত্তর তৈরি করে।
\item দৃষ্টিভিত্তিক মডেল পিক্সেল বা ভিডিওর দৃশ্য থেকে শ্রেণি, অবস্থান বা বর্ণনা নির্ণয় করে। CNN-ভিত্তিক শ্রেণিনির্ণয় দিয়ে পাতার রোগ বা পণ্যের ত্রুটি শনাক্ত করা যায়।
\item ছবি তৈরির মডেল লেখা বা আগের ছবি থেকে নতুন বা সম্পাদিত ছবি দেয়। GPT Image এবং Latent Diffusion-ভিত্তিক মডেল উদাহরণ।
\item শব্দের মডেলে অডিও থেকে লেখা অথবা লেখা থেকে অডিও পাওয়া যায়। Whisper কথা থেকে লেখা তৈরির উদাহরণ; পাঠ শোনানোর জন্য আলাদা Text-to-Speech মডেল লাগতে পারে।
\item Embedding মডেল লেখা বা সমর্থিত অন্য ডেটাকে সংখ্যার ভেক্টরে প্রকাশ করে। অর্থের মিল ধরে নোট খুঁজতে এটি লাগে; ভেক্টর নিজে প্রস্তুত উত্তর নয়।
\item পূর্বাভাসের মডেল রেকর্ড থেকে শ্রেণি, সম্ভাবনা বা পরিমাণ নির্ণয় করে। সব বিক্রির পূর্বাভাস বা স্প্যাম শনাক্তকারী মডেল LLM নয়।
\item Multimodal মডেল একাধিক ধরনের ইনপুটের সম্পর্ক বিশ্লেষণ করতে পারে; ফলের ধরন নির্দিষ্ট মডেলের ওপর নির্ভর করে।
\end{compactpoints}
এই শ্রেণিগুলো পরস্পর সম্পূর্ণ পৃথক নয়। Transformer, CNN ও Diffusion গঠন বা পদ্ধতির নাম; GPT, Claude ও Gemini মডেল পরিবারের নাম। প্রচলিত Denoising Diffusion-এ প্রশিক্ষণের সময় নয়েজ ও ছবির সম্পর্ক শেখানো হয়। তৈরির সময় নয়েজ থেকে ধাপে ধাপে কাঙ্ক্ষিত গঠন পাওয়া যায়; প্রম্পট দিকনির্দেশ দিতে পারে। এটি ভাষা মডেলের পরবর্তী টোকেন নির্বাচনের একই প্রক্রিয়া নয়।
\term{Open Weight}-এ শেখানো প্যারামিটার ব্যবহারের সুযোগ থাকে। প্রশিক্ষণ-ডেটা, কোড ও সব ব্যবহারের অধিকার প্রকাশিত নাও থাকতে পারে; তাই একে স্বয়ংক্রিয়ভাবে সম্পূর্ণ Open Source বলা যায় না। লাইসেন্স, যন্ত্র ও ব্যবহারের শর্ত দেখতে হয়।
\subtopichead{উপযুক্ত মডেল বাছাই}
কলেজের নোট-সহায়কে স্ক্যানের জন্য OCR, প্রাসঙ্গিক নোট খুঁজতে Embedding ও Retrieval, ব্যাখ্যায় LLM এবং পাঠ শোনাতে শব্দের মডেল লাগতে পারে। সব কাজ এক মডেলে করতে হয় না। শিক্ষক, মূল বই ও যাচাই করা ফল দিয়ে উপযোগিতা পরীক্ষা করতে হয়।
মডেল বাছাইয়ে কাজ, ডেটা, ভাষা, গতি, ব্যয় ও ভুলের মাপকাঠি বিবেচনা করা হয়। বড় মডেলও ভুল করতে পারে; ছোট মডেল নির্দিষ্ট কাজ দ্রুত করতে পারে। কিছু Reasoning মডেল জটিল উত্তরের জন্য বেশি গণনা বা বহু ধাপের বিশ্লেষণ করে; সময় ও ব্যয় বাড়লেও নির্ভুলতা নিশ্চিত হয় না।
নিজস্ব কম্পিউটারে মডেল চালালে উপযুক্ত হার্ডওয়্যার ও পরিচালনা দরকার। Cloud বা API দিয়ে ব্যবহার করলে সংযোগ, ব্যয় ও ডেটা ব্যবহারের শর্ত বিবেচনা করতে হয়। নাম, প্যারামিটার সংখ্যা বা প্রচারের ক্রমের বদলে একই ধরনের বাস্তব সমস্যা দিয়ে ফল পরীক্ষা করা দরকার।
